% ===================================================================
%  ТАБЛИЦА № 11 — Город и его здания. — Пожар.
%  Traduction 2026 d'apres texto/io, controlee sur texto/fr, texto/en
%  et texto/es. Consigne : voir texto/ru/10-tablica-01.tex.
% ===================================================================

%%K t11-apar-1 apar
\VUsaut{2.0mm}
\VUcentre{13.2pt}{90}{ТРЕТЬЯ СЕРИЯ}
\VUsaut{3.0mm}
\VUfilet{16mm}
\VUsaut{5.6mm}
\VUtitre{15.0pt}{60}{ТАБЛИЦА № 11}
\VUsaut{1.4mm}
\VUpk{11.4pt}{40}{Город и его здания. --- Пожар.}
\VUsaut{2.2mm}

%%K t11-tit-1 sub
\VUpk{10.2pt}{40}{Окраины.}
\VUsaut{3.0mm}

%%K t11-c1-01-1 p
1. --- Я живу в большом городе в ста километрах от океана, который тем не
менее является первоклассным \VUgras{портом} \textsuperscript{(1)}. Река, что течёт
вдоль него, шириной в четыреста метров и образует прекрасный рейд.

%%K t11-c1-02-1 p
2. --- Вся часть города вдоль \VUgras{набережных} \textsuperscript{(2)} имеет вид
совершенно морской и промышленный и образует \VUgras{предместье}
\textsuperscript{(3)}, где живут одни моряки да рабочие. Последние работают на
\VUgras{заводах} \textsuperscript{(4)} и на \VUgras{газовом заводе} \textsuperscript{(5)},
чью высокую \VUgras{трубу} \textsuperscript{(6)} и \VUgras{газгольдер} видно
издалека.

%%K t11-c1-03-1 p
3. --- В Средние века город был укреплён, и кое-какие остатки его валов
ещё находят, в особенности \VUgras{ворота} \textsuperscript{(8)} старого
\VUgras{укреплённого замка} \textsuperscript{(7)} с двумя \VUgras{башнями}
\textsuperscript{(9)} по бокам, которые теперь служат \VUgras{тюрьмой}.

%%K t11-c1-04-1 p
4. --- Во всём этом \VUgras{квартале}, на вид очень старом, только одно
здание сравнительно новое: \VUgras{таможня} \textsuperscript{(10)}. Она стоит
между набережной и маленькой \VUgras{улицей} \textsuperscript{(11)}, ведущей к
старой \VUgras{ратуше} \textsuperscript{(12)}. Это последнее здание легко узнать
по исполинской \VUgras{дозорной башне} \textsuperscript{(13)}, которая возвышается
над старым городом.

%%K t11-c1-05-1 p
5. --- По другую сторону улицы \VUgras{стоит} здание в том же стиле, что и
таможня: это \VUgras{биржа} \textsuperscript{(14)}. Один из её фасадов выходит на
маленькую площадь, где находится \VUgras{префектура} \textsuperscript{(15)}. Наш
город, правда, не столица, но всё же важный уездный центр.

%%K t11-tit-2 sub
\VUsaut{5.0mm}
\VUpk{10.2pt}{40}{Верхняя часть города.}
\VUsaut{3.4mm}

%%K t11-c1-06-1 p
6. --- Гарнизон, довольно большой, помещается в обширных и очень здоровых
\VUgras{казармах} \textsuperscript{(16)}; они тянутся до самой решётки
\VUgras{городского парка} \textsuperscript{(17)}. \VUgras{Бульвар} \textsuperscript{(19)},
ведущий к этому парку, проходит позади \VUgras{сберегательной кассы}
\textsuperscript{(18)} и выводит на площадь \VUgras{собора} \textsuperscript{(20)}.

%%K t11-c1-07-1 p
7. --- Все эти здания подымаются уступами на небольшом холме, где стоит и
наш большой \VUgras{лицей} \textsuperscript{(21)}.

%%K t11-c1-07-2 p
Если спуститься по слегка покатой улице, выходящей на Большую улицу,
дойдёшь до \VUgras{суда} \textsuperscript{(22)}, перед которым лежит приятный
\VUgras{сквер} \textsuperscript{(26)}. \VUgras{Площадь} эта невелика, но посреди
города она образует оазис зелени и прохлады. Поэтому её охотно посещают
горожане, которые любят посидеть под навесом \VUgras{кафе} \textsuperscript{(25)},
слушая военный оркестр, что играет по четвергам и воскресеньям в
\VUgras{беседке} \textsuperscript{(27)} среди газонов и клумб.

%%K t11-c1-08-1 p
8. --- На одном из этих газонов стоит бронзовая \VUgras{статуя}
\textsuperscript{(28)}, памятник, поставленный в память одного из наших горожан,
оказавшего большие услуги родному городу.

%%K t11-tit-3 sub
\VUsaut{4.0mm}
\VUpk{10.2pt}{40}{Передвижение.}
\VUsaut{3.4mm}

%%K t11-c1-09-1 p
9. --- Наш город ещё не освещён электричеством; у него только
\VUgras{газовые фонари} \textsuperscript{(29)}. Но \VUgras{местная печать} хлопочет
о том, чтобы этот способ освещения был улучшен, \VUgras{так же как} уже
улучшена \VUgras{тяга} на трамваях. \VUgras{Старые} вагоны, \VUgras{которые
тянули} \VUgras{лошади}, заменены удобными \VUgras{электрическими трамваями}
\textsuperscript{(31)}, которые быстро возят пассажиров из одного конца
\VUgras{города в другой} за очень \VUgras{низкую} плату.

%%K t11-c1-10-1 p
10. --- Возле суда --- \VUgras{банк} \textsuperscript{(32)}, где учитывают торговые
векселя. \VUgras{Банковские артельщики} \textsuperscript{(33)} ходят по домам
собирать причитающееся.

%%K t11-tit-4 sub
\VUsaut{4.0mm}
\VUpk{10.2pt}{40}{Искусство и наука.}
\VUsaut{3.4mm}

%%K t11-c1-11-1 p
11. --- Красивое здание неподалёку --- \VUgras{музей}. В нём помещаются
разом городская \VUgras{библиотека} \textsuperscript{(34)}, прекрасные
\VUgras{естественнонаучные собрания} \textsuperscript{(35)} (естественноисторический
музей) и изящная \VUgras{картинная галерея} \textsuperscript{(36)}, составляющая
великолепную постоянную \VUgras{выставку}.

%%K t11-c1-11-2 p
Он открыт для публики два раза в неделю, но посетители могут осмотреть
его в любой день за небольшую подачку \VUgras{сторожу} \textsuperscript{(37)}.
\VUgras{Художники} \textsuperscript{(38)} приходят туда работать. Их видишь
\VUgras{перед} мольбертами, с \VUgras{палитрой} в руке, копирующими знаменитые
картины.

%%K t11-c1-12-1 p
12. --- На высотах за музеем едва различаешь \VUgras{купол} \textsuperscript{(40)}
\VUgras{больницы} \textsuperscript{(39)} и здания \VUgras{учительской семинарии}
\textsuperscript{(41)}, где учились преподаватели наших городских школ. На
Театральной площади есть \VUgras{почта} \textsuperscript{(43)}, помещённая в
\VUgras{частном доме} \textsuperscript{(42)}.

%%K t11-tit-5 sub
\VUsaut{5.0mm}
\VUpk{11.4pt}{40}{Пожар.}
\VUsaut{3.0mm}

%%K t11-tit-6 sub
\VUpk{10.2pt}{40}{Крик «Пожар!».}
\VUsaut{3.4mm}

%%K t11-c2-01-1 p
1. --- По городу разносится тревожная весть: в театре только что начался
пожар! Когда я добираюсь до \VUgras{площади} \textsuperscript{(96)}, толпа уже
собралась на \VUgras{тротуаре} \textsuperscript{(46)} и на \VUgras{мостовой}
\textsuperscript{(47)}. \VUgras{Почтовые чиновники} \textsuperscript{(44)} и
\VUgras{почтальоны} \textsuperscript{(45)} выходят с почты. \VUgras{Вагоновожатый}
\textsuperscript{(48)} принуждён был остановить свой трамвай, чтобы пропустить
городскую \VUgras{карету скорой помощи} \textsuperscript{(50)}, прибывшую с
\VUgras{хирургом} \textsuperscript{(52)} и \VUgras{сиделкой} \textsuperscript{(51)},
чтобы помочь \VUgras{раненому} \textsuperscript{(53)}, которого на \VUgras{носилках}
\textsuperscript{(54)} несут к месту возле \VUgras{газетного киоска}
\textsuperscript{(55)}.

%%K t11-c2-02-1 p
2. --- Рота пехоты, возвращавшаяся с ученья, остановилась, чтобы помочь
поддерживать порядок вместе с конными \VUgras{городскими стражниками}
\textsuperscript{(61)}. \VUgras{Всадники}, чьи лошади встают на дыбы, справляются
лучше \VUgras{солдат} \textsuperscript{(57)} и \VUgras{жандармов}
\textsuperscript{(63)}, оттесняя \VUgras{зевак} \textsuperscript{(56)}. Жандармы к тому же
очень заняты. Один из них указывает \VUgras{полицейским} \textsuperscript{(64)},
куда нести другого \VUgras{пострадавшего} \textsuperscript{(65)}, храброго
пожарного, упавшего с лестницы.

%%K t11-c2-03-1 p
3. --- \VUgras{Офицер} \textsuperscript{(66)} указывает в точности, куда поставить
\VUgras{паровую помпу} \textsuperscript{(67)}, которая как раз подъезжает. В
ожидании этой мощной машины он велел поставить \VUgras{ручную помпу}
\textsuperscript{(68)}, длинный \VUgras{рукав} \textsuperscript{(69)} которой льёт на
огонь потоки воды. Её качают шестеро пожарных, а их товарищ, держа
\VUgras{брандспойт} \textsuperscript{(70)}, направляет \VUgras{струю}
\textsuperscript{(71)} в самое сердце пожара.

%%K t11-c2-04-1 p
4. --- По другую сторону театра подняли большую \VUgras{лестницу}
\textsuperscript{(72)}. Она позволяет пожарным подобраться к пламени, что
взвивается среди клубов чёрного \VUgras{дыма} \textsuperscript{(75)}.

%%K t11-tit-7 sub
\VUsaut{5.0mm}
\VUpk{10.2pt}{40}{Смелые дела.}
\VUsaut{3.4mm}

%%K t11-c2-05-1 p
5. --- \VUgras{Огонь} \textsuperscript{(74)} начался в фойе \VUgras{театра}
\textsuperscript{(73)}, во время репетиции \VUgras{оперы}, первое
\VUgras{представление} которой должно было состояться завтра. К счастью, в
зале было лишь несколько \VUgras{зрителей} \textsuperscript{(76)}. Бедные люди
укрылись на \VUgras{балконе} \textsuperscript{(77)}, куда храбрые \VUgras{пожарные}
\textsuperscript{(86)} идут им на помощь с лестницей и верёвками.

%%K t11-c2-06-1 p
6. --- Под \VUgras{портиком} \textsuperscript{(78)}, где \VUgras{афиша}
\textsuperscript{(79)} объявляет о представлении, видно, как бегут
\VUgras{музыканты} \textsuperscript{(81)} \VUgras{оркестра} \textsuperscript{(80)},
унося свои инструменты: \VUgras{скрипку} \textsuperscript{(81)}, \VUgras{флейту}
\textsuperscript{(82)}, \VUgras{виолончель} \textsuperscript{(83)}, \VUgras{тромбон}
\textsuperscript{(84)}, \VUgras{барабан} \textsuperscript{(85)} и так далее. Стараются
спасти часть \VUgras{обстановки} \textsuperscript{(87)}.

%%K t11-c2-07-1 p
7. --- \VUgras{Актёрам} \textsuperscript{(89)} и \VUgras{актрисам} \textsuperscript{(88)},
\VUgras{певцам} и певицам, пришлось выбежать в сценических \VUgras{костюмах}
\textsuperscript{(90)}. Тенор объясняет \VUgras{журналисту} \textsuperscript{(92)}, как
это случилось. \VUgras{Репортёр} прибежал с самой первой минуты вместе с
властями: \VUgras{префектом} \textsuperscript{(91)}, \VUgras{мэром}
\textsuperscript{(93)}, \VUgras{полицмейстером} \textsuperscript{(94)} и
\VUgras{следователем} \textsuperscript{(95)}, который поведёт дознание, чтобы
установить, на ком лежит ответственность.
\VUsaut{6.0mm}
\VUfilet{22mm}
